\chapter{The recorder}
\label{ch:u-recorder}

The robot records its cameras on its own disk: the visible and thermal video, every raw thermal
frame (temperatures per pixel, for analysis later) with the robot's position, and the acoustic
camera's data. Nothing goes over the network while it records, so a weak link loses nothing.

\section{When it records}

\begin{itemize}
  \item \textbf{Every patrol}, by itself: from the moment the robot walks onto the route until
    30~s after it finished.
  \item \textbf{On request}: start and stop it from the H12 (the recorder action) or in NETRA
    (\menupath{Live control › Recorder}, the switch). A manual recording that runs into a patrol
    goes on as one session.
\end{itemize}

\section{What you see}

\begin{itemize}
  \item \textbf{H12}: \ui{REC 4:12 · 250 frames · patrol}, or \ui{Recorder ready}; \ui{Recorder
    not answering} when the perception host does not answer.
  \item \textbf{NETRA}: a red \ui{REC} badge on the robot's card; in \menupath{Live control} the
    duration, who started it, the thermal frames, the videos, the free disk, and a warning when
    frames are dropped because the disk is too slow.
\end{itemize}

\section{The disk}

Before each session and every 30~s while it runs, the oldest recordings are removed until 10~GB
is free; below 3~GB nothing is recorded and the recorder says why. Copy recordings you want to
keep off the robot (the technician: \file{~/x30/recordings} on the perception host).
